\section{Information conventions}
\label{sec:v32-ledger}
The following conventions specify the information and resource used in
individual results. A clock or a fixed program does not contain observed
data. A public descriptor or a persistent seed is supplied only where
expressly stated. Separate rows below are not assertions of equivalence
between decision problems.

\paragraph{The collision classification and its finite certificates.}
The acquisition commands are prescribed, independent and uniform on the
command cube. The latent parameter is shared by all trials. The update
reads the previous label, the current command and the current report;
memory is charged after every report. The decoder reads only the label,
the clock and a query selected independently after compression. Past
commands, reports, predictions and retained seeds are unavailable.
Private update coins are fresh and discarded. One selected checkpoint and
its one remaining continuation are executed in a run; scores do not feed
later updates. The target is a fixed spanning menu of the entire relevant
future experiment, not only a newly acquired current block. The two risks
are
\[
 \inf_F\max_t\mathbb E_h\mathbb E_{\rm code} L_t,
 \qquad
 \inf_F\max_t\sup_h\mathbb E_{\rm code} L_t.
\]
The finite approximation and relative certificate use this same resource.
The calibration, prior and horizon are fixed; constants are uniform in
calibration within its prescribed compact chamber and in the label budget.

\paragraph{State-controlled acquisition.}
In the additional controlled model, the next command is drawn from the
preceding retained label before the new report. There is no free
full-history scheduler. The decoder, charge boundary, shared latent
parameter and exclusion of a persistent seed are as in the preceding
paragraph. The objective is $\inf_F\max_t\mathbb E_F L_t$; its law is
induced by the common controller. The compatibility and command-rounding
results apply. The exploration-law collision lower bound is not claimed
for acquisition optimized jointly with the controller.

\paragraph{Independent-edge graph experiments.}
The inherited fixed-graph model charges memory at vertex-batch boundaries
and uses full tensor-product future probes. Its adaptive converse and
resolution-blind acquisition comparisons retain the information and
finite descriptors explicitly permitted in their definitions; that
converse may be a relaxation of the implementing scheduler. The
size-uniform result uses a different, local-and-tensor query weighting
under its retained balanced-cut hypotheses. Equality of test spans is not
uniform equivalence of these loss metrics as graph size grows. Neither
the new finite certificate nor the regenerative bound changes that task.

\paragraph{Committed and predictable regenerative networks.}
The route or current vertex is an explicitly public finite descriptor,
and an independent persistent seed is available to the decoder as well
as to the controller. At a charged block boundary the decoder also
receives one of $M$ labels and the current query. A committed route is
selected before data. In the predictable network converse the scheduler
may read the revealed history, including earlier query outcomes, but
selects a vertex before observing its block. Conditional on that history,
the seed and the selected vertex, the new block has its specified law
and is sufficient for the current query. The minimizing construction is
an open-loop randomized path and reset encoding. The maximum of expected
checkpoint losses is taken after seed averaging. These statements are
not an equivalence with the unseeded model, and a charged descriptor is a
different resource with the stated additional-state upper bound.

\paragraph{The positive delayed-use comparison.}
The two deterministic schedules are announced in advance. Each acquires
the same four reports once and executes one selected checkpoint query.
Memory is charged after each acquisition block and at both checkpoints;
the complete block may be read during its update. A decoder reads the
label and clock, without a retained public seed or past-output access.
The selected score is not fed back. The objective is the maximum of the
two expected current-query excesses, minimized over one controller.
Just-in-time and early acquisition are both admitted schedules of this
one timing problem. Their common label and seed convention agrees, but
their block-boundary convention is not silently identified with the
per-report convention of the monomial theorem.
